\chapter{Annual schedule and time slots}\label{ANNUEL}
70 occurrences present; 2 slots neutralised for the mock exam windows; 4 slots after 18 June conditional. 64 ordinary sessions before 18 June, i.e. 128 h, including two full practice papers. Saturday tests and mock exams: duration and placement to be confirmed, not artificially added to the teaching volume.
\paragraph{120-minute structure.} Fluency routine 8; homework review 10; investigation or revisiting 15; lesson and example 30; exercises 40; wrap-up and homework 7; flexibility margin 10. For a first session, the diagnostic takes part of the exercise time. For a full paper, the 120 min are reserved for the paper, with no fluency routine added.
\Needspace{8\baselineskip}\section*{S01 -- 2026-09-03: Welcome and cross-curricular diagnostic}
\textbf{Status:} lesson. 09:45–11:45.\par\textbf{Fluency routine:} R001, R008.\par\textbf{In class:} Diagnostic of Year 10 (seconde) prior knowledge and presentation of working methods.\par\textbf{Homework:} Finish only the diagnostic questions not yet covered, 10 min.\par\textbf{Note:} Expected mastery, to be adjusted from the work actually done and the diagnostics.
\Needspace{8\baselineskip}\section*{S02 -- 2026-09-08: Defining and representing a sequence}
\textbf{Status:} lesson. 09:45–11:45.\par\textbf{Fluency routine:} R002, R009.\par\textbf{In class:} C01-U1; worked examples and applications; C01-E1. Diagnostic D: 10 min; investigation RECH: 15 min.\par\textbf{Homework:} CQFD, exercise 9 p. 64, parts a and b. Calculate the first four terms. Write out the substitutions before simplifying. 10–20 min maximum; replaces the initial homework.\par\textbf{Note:} Expected mastery, to be adjusted from the work actually done and the diagnostics. CQFD solutions in the teacher booklet; return to it at the next lesson.
\Needspace{8\baselineskip}\section*{S03 -- 2026-09-10: Arithmetic sequences}
\textbf{Status:} lesson. 09:45–11:45.\par\textbf{Fluency routine:} R003, R031.\par\textbf{In class:} C01-U2; worked examples and applications; C01-E2: version A only. \par\textbf{Homework:} CQFD, exercise 11 p. 64, all parts. Calculate the terms successively and keep a record of the four steps. 10–20 min maximum; replaces the initial homework.\par\textbf{Note:} Expected mastery, to be adjusted from the work actually done and the diagnostics. CQFD solutions in the teacher booklet; return to it at the next lesson.
\Needspace{8\baselineskip}\section*{S04 -- 2026-09-15: Geometric sequences}
\textbf{Status:} lesson. 09:45–11:45.\par\textbf{Fluency routine:} R004, R032.\par\textbf{In class:} C01-U3; worked examples and applications; C01-E3; E2 version B. \par\textbf{Homework:} C01-E3, 10–20 min maximum.\par\textbf{Note:} Expected mastery, to be adjusted from the work actually done and the diagnostics.
\Needspace{8\baselineskip}\section*{S05 -- 2026-09-17: Sums, indices and proofs}
\textbf{Status:} lesson. 09:45–11:45.\par\textbf{Fluency routine:} R005, R032.\par\textbf{In class:} C01-U4, C01-U5; worked examples and applications; C01-E4; E6 as time allows.  MCQ 1 C01 (10 min, including correction). MCQ 2 C01 (10 min) and wrap-up.\par\textbf{Homework:} C01-E8 to revisit in a week, 10–20 min maximum.\par\textbf{Note:} Expected mastery, to be adjusted from the work actually done and the diagnostics. Sums: essential formulas and proofs; exploration of thresholds to be revisited on 15 December. No formal induction.
\Needspace{8\baselineskip}\section*{S06 -- 2026-09-22: Three forms, three readings}
\textbf{Status:} lesson. 09:45–11:45.\par\textbf{Fluency routine:} R006, R031.\par\textbf{In class:} C02-U1; worked examples and applications; C02-E1. Diagnostic D: 10 min; investigation RECH: 15 min.\par\textbf{Homework:} C02-E1, 10–20 min maximum.\par\textbf{Note:} Expected mastery, to be adjusted from the work actually done and the diagnostics.
\Needspace{8\baselineskip}\section*{S07 -- 2026-09-24: Discriminant and solving}
\textbf{Status:} lesson. 09:45–11:45.\par\textbf{Fluency routine:} R007, R031.\par\textbf{In class:} C02-U2; worked examples and applications; C02-E3. \par\textbf{Homework:} CQFD, exercise 43 p. 99, parts a and c. First find an obvious root, then the other using the product of the roots. 10–20 min maximum; replaces the initial homework.\par\textbf{Note:} Expected mastery, to be adjusted from the work actually done and the diagnostics. CQFD solutions in the teacher booklet; return to it at the next lesson.
\Needspace{8\baselineskip}\section*{S08 -- 2026-09-29: Factorising without automating the discriminant}
\textbf{Status:} lesson. 09:45–11:45.\par\textbf{Fluency routine:} R008, R031.\par\textbf{In class:} C02-U3; worked examples and applications; C02-E4. \par\textbf{Homework:} CQFD, exercise 49 p. 99, parts a and b. Translate the sum and the product into an equation; do not assume in advance that the numbers exist. 10–20 min maximum; replaces the initial homework.\par\textbf{Note:} Expected mastery, to be adjusted from the work actually done and the diagnostics. CQFD solutions in the teacher booklet; return to it at the next lesson.
\Needspace{8\baselineskip}\section*{S09 -- 2026-10-01: Sign, inequalities and optimisation}
\textbf{Status:} lesson. 09:45–11:45.\par\textbf{Fluency routine:} R009, R031.\par\textbf{In class:} C02-U4; worked examples and applications; C02-E5 or E6.  MCQ 1 C02 (10 min, including correction). MCQ 2 C02 (10 min) and wrap-up.\par\textbf{Homework:} C02-E8 to revisit in a week, 10–20 min maximum.\par\textbf{Note:} Expected mastery, to be adjusted from the work actually done and the diagnostics.
\Needspace{8\baselineskip}\section*{S10 -- 2026-10-06: Projection and angle}
\textbf{Status:} lesson. 09:45–11:45.\par\textbf{Fluency routine:} R010, R031.\par\textbf{In class:} C03-U1; worked examples and applications; C03-E4 and E6. Diagnostic D: 10 min; investigation RECH: 15 min.\par\textbf{Homework:} C03-E4, 10–20 min maximum.\par\textbf{Note:} Expected mastery, to be adjusted from the work actually done and the diagnostics.
\Needspace{8\baselineskip}\section*{S11 -- 2026-10-08: Coordinates, norm and orthogonality}
\textbf{Status:} lesson. 09:45–11:45.\par\textbf{Fluency routine:} R011, R031.\par\textbf{In class:} C03-U2; worked examples and applications; C03-E1 and E2. \par\textbf{Homework:} CQFD, exercise 9 p. 250, parts a and b. Calculate the dot product in the orthonormal basis and detail the two products of coordinates. 10–20 min maximum; replaces the initial homework.\par\textbf{Note:} Expected mastery, to be adjusted from the work actually done and the diagnostics. CQFD solutions in the teacher booklet; return to it at the next lesson.
\Needspace{8\baselineskip}\section*{S12 -- 2026-10-13: Vector calculation and identities}
\textbf{Status:} lesson. 09:45–11:45.\par\textbf{Fluency routine:} R012, R031.\par\textbf{In class:} C03-U3; worked examples and applications; C03-E3.  MCQ 1 C03 (10 min, including correction).\par\textbf{Homework:} CQFD, exercise 11 p. 250, questions 1, 2 and 3. Use in turn the coordinates, Chasles' relation and the fact that I is the midpoint. 10–20 min maximum; replaces the initial homework.\par\textbf{Note:} Expected mastery, to be adjusted from the work actually done and the diagnostics. CQFD solutions in the teacher booklet; return to it at the next lesson.
\Needspace{8\baselineskip}\section*{S13 -- 2026-10-15: Al-Kashi and choosing a method}
\textbf{Status:} lesson. 09:45–11:45.\par\textbf{Fluency routine:} R013, R031.\par\textbf{In class:} C03-U4, C03-U5; worked examples and applications; C03-E8 and E5 as time allows.  MCQ 2 C03 (10 min) and wrap-up.\par\textbf{Homework:} C03-E8 to revisit in a week, 10–20 min maximum.\par\textbf{Note:} Expected mastery, to be adjusted from the work actually done and the diagnostics.
\Needspace{8\baselineskip}\section*{S14 -- 2026-11-03: Rate of change and secants}
\textbf{Status:} lesson. 09:45–11:45.\par\textbf{Fluency routine:} R014, R038.\par\textbf{In class:} C04-U1; worked examples and applications; C04-E1: compute the two rates only; comparison with the derivative at a point deferred. Diagnostic D: 10 min; investigation RECH: 15 min. Made available: EV1-T1: choose a question suited to the topics already studied.\par\textbf{Homework:} C04-E1: the rates only, 10–20 min maximum.\par\textbf{Note:} Expected mastery, to be adjusted from the work actually done and the diagnostics. Spare practice paper: replaces a planned homework task, with no accumulation of homework; do not require two full papers at home.
\Needspace{8\baselineskip}\section*{S15 -- 2026-11-05: Derivative at a point}
\textbf{Status:} lesson. 09:45–11:45.\par\textbf{Fluency routine:} R015, R039.\par\textbf{In class:} C04-U2; worked examples and applications; C04-E2 and E3.  Made available: EV1-T2: choose a question suited to the topics already studied.\par\textbf{Homework:} CQFD, exercise 17 p. 127, all parts. Read two points on each tangent; write the quotient of the increments rather than estimating its slope by eye. 10–20 min maximum; replaces the initial homework.\par\textbf{Note:} Expected mastery, to be adjusted from the work actually done and the diagnostics. Spare practice paper: replaces a planned homework task, with no accumulation of homework; do not require two full papers at home. CQFD solutions in the teacher booklet; return to it at the next lesson.
\Needspace{8\baselineskip}\section*{S16 -- 2026-11-10: Tangent and linear approximation}
\textbf{Status:} lesson. 09:45–11:45.\par\textbf{Fluency routine:} R016, R040.\par\textbf{In class:} C04-U3; worked examples and applications; C04-E4 and E5.  MCQ 1 C04 (10 min, including correction).\par\textbf{Homework:} CQFD, exercise 19 p. 127, questions 1 and 2. Distinguish the y-coordinate of the point of contact from the slope of the tangent line. 10–20 min maximum; replaces the initial homework.\par\textbf{Note:} Expected mastery, to be adjusted from the work actually done and the diagnostics. CQFD solutions in the teacher booklet; return to it at the next lesson.
\Needspace{8\baselineskip}\section*{S17 -- 2026-11-12: When differentiability fails}
\textbf{Status:} lesson. 09:45–11:45.\par\textbf{Fluency routine:} R017, R041.\par\textbf{In class:} C04-U4; worked examples and applications; C04-AP4 then E7 or E8.  MCQ 2 C04 (10 min) and wrap-up.\par\textbf{Homework:} C04-E8 to revisit in a week, 10–20 min maximum.\par\textbf{Note:} Expected mastery, to be adjusted from the work actually done and the diagnostics.
\Needspace{8\baselineskip}\section*{S18 -- 2026-11-17: Conditioning: changing the denominator}
\textbf{Status:} lesson. 09:45–11:45.\par\textbf{Fluency routine:} R018, R042.\par\textbf{In class:} C05-U1; worked examples and applications; C05-E1. Diagnostic D: 10 min; investigation RECH: 15 min.\par\textbf{Homework:} C05-E1, 10–20 min maximum.\par\textbf{Note:} Expected mastery, to be adjusted from the work actually done and the diagnostics.
\Needspace{8\baselineskip}\section*{S19 -- 2026-11-19: Trees, partitions and total probability}
\textbf{Status:} lesson. 09:45–11:45.\par\textbf{Fluency routine:} R019, R043.\par\textbf{In class:} C05-U2; worked examples and applications; C05-E2 and E6. \par\textbf{Homework:} CQFD, exercise 67 p. 336, questions 1, 2 and 3. Build a tree keeping the events D and N from the textbook. Note: N denotes a negative test. 10–20 min maximum; replaces the initial homework.\par\textbf{Note:} Expected mastery, to be adjusted from the work actually done and the diagnostics. CQFD solutions in the teacher booklet; return to it at the next lesson.
\Needspace{8\baselineskip}\section*{S20 -- 2026-11-24: Independence and incompatibility}
\textbf{Status:} lesson. 09:45–11:45.\par\textbf{Fluency routine:} R020, R044.\par\textbf{In class:} C05-U3; worked examples and applications; C05-E3.  MCQ 1 C05 (10 min, including correction).\par\textbf{Homework:} C05-E3, 10–20 min maximum.\par\textbf{Note:} Expected mastery, to be adjusted from the work actually done and the diagnostics.
\Needspace{8\baselineskip}\section*{S21 -- 2026-11-26: Successive trials and Bernoulli}
\textbf{Status:} lesson. 09:45–11:45.\par\textbf{Fluency routine:} R021, R045.\par\textbf{In class:} C05-U4, C05-U5; worked examples and applications; C05-E4 and E5 as preferred.  MCQ 2 C05 (10 min) and wrap-up.\par\textbf{Homework:} CQFD, exercise 70 p. 336, parts a and b. Count the four equally likely outcomes of the two throws to refute the proposed equalities. 10–20 min maximum; replaces the initial homework.\par\textbf{Note:} Expected mastery, to be adjusted from the work actually done and the diagnostics. CQFD solutions in the teacher booklet; return to it at the next lesson.
\Needspace{8\baselineskip}\section*{S22 -- 2026-12-01: Derivative function and standard functions}
\textbf{Status:} lesson. 09:45–11:45.\par\textbf{Fluency routine:} R022, R040.\par\textbf{In class:} C06-U1; worked examples and applications; C06-AP1: standard derivatives only. Diagnostic D: 10 min; investigation RECH: 15 min.\par\textbf{Homework:} C06-AP1, 10–20 min maximum.\par\textbf{Note:} Expected mastery, to be adjusted from the work actually done and the diagnostics.
\Needspace{8\baselineskip}\section*{S23 -- 2026-12-03: Sum, product, reciprocal and quotient}
\textbf{Status:} lesson. 09:45–11:45.\par\textbf{Fluency routine:} R023, R041.\par\textbf{In class:} C06-U2; worked examples and applications; C06-E1. \par\textbf{Homework:} C06-E1, 10–20 min maximum.\par\textbf{Note:} Expected mastery, to be adjusted from the work actually done and the diagnostics.
\Needspace{8\baselineskip}\section*{S24 -- 2026-12-08: Sign of the derivative and extrema}
\textbf{Status:} lesson. 09:45–11:45.\par\textbf{Fluency routine:} R024, R042.\par\textbf{In class:} C06-U3; worked examples and applications; C06-E2.  MCQ 1 C06 (10 min, including correction).\par\textbf{Homework:} CQFD, exercise 16 p. 156, parts f and g. For each function, give an interval on which it is both negative and increasing; justify the two properties separately. 10–20 min maximum; replaces the initial homework.\par\textbf{Note:} Expected mastery, to be adjusted from the work actually done and the diagnostics. CQFD solutions in the teacher booklet; return to it at the next lesson.
\Needspace{8\baselineskip}\section*{S25 -- 2026-12-10: Parity, inequalities and comparison}
\textbf{Status:} lesson. 09:45–11:45.\par\textbf{Fluency routine:} R025, R043.\par\textbf{In class:} C06-U4, C06-U5; worked examples and applications; C06-E4 and E5 as preferred.  MCQ 2 C06 (10 min) and wrap-up.\par\textbf{Homework:} CQFD, exercise 12 p. 156, questions 1 and 2. Calculate the difference, study its sign and translate this into the relative position of the curves. 10–20 min maximum; replaces the initial homework.\par\textbf{Note:} Expected mastery, to be adjusted from the work actually done and the diagnostics. CQFD solutions in the teacher booklet; return to it at the next lesson.
\Needspace{8\baselineskip}\section*{S26 -- 2026-12-15: Revisiting sums and thresholds}
\textbf{Status:} lesson. 09:45–11:45.\par\textbf{Fluency routine:} R026, R038.\par\textbf{In class:} C01-E4, E6, E7: check the number of terms and the preceding index; revisit the dot product C03-E8. Made available: EV2-T1: choose a question suited to the topics already studied.\par\textbf{Homework:} C01-E8, 10 min.\par\textbf{Note:} Expected mastery, to be adjusted from the work actually done and the diagnostics. Spare practice paper: replaces a planned homework task, with no accumulation of homework; do not require two full papers at home.
\Needspace{8\baselineskip}\section*{S27 -- 2026-12-17: Consolidating differentiation and probability}
\textbf{Status:} lesson. 09:45–11:45.\par\textbf{Fluency routine:} R027, R045.\par\textbf{In class:} C06-E1 and E2; compare with C05-E8, two separate consolidation phases. Made available: EV2-T2: choose a question suited to the topics already studied.\par\textbf{Homework:} Review the errors identified, 15 min.\par\textbf{Note:} Expected mastery, to be adjusted from the work actually done and the diagnostics. Spare practice paper: replaces a planned homework task, with no accumulation of homework; do not require two full papers at home.
\Needspace{8\baselineskip}\section*{S28 -- 2027-01-05: Normal vector and Cartesian equation}
\textbf{Status:} lesson. 09:45–11:45.\par\textbf{Fluency routine:} R028, R052.\par\textbf{In class:} C07-U1; worked examples and applications; C07-E1. Diagnostic D: 10 min; investigation RECH: 15 min.\par\textbf{Homework:} CQFD, exercise 18 p. 279, parts a, b, c and d. Read a non-zero normal vector for each equation, without solving for y. 10–20 min maximum; replaces the initial homework.\par\textbf{Note:} Expected mastery, to be adjusted from the work actually done and the diagnostics. CQFD solutions in the teacher booklet; return to it at the next lesson.
\Needspace{8\baselineskip}\section*{S29 -- 2027-01-07: Orthogonal projection and distance}
\textbf{Status:} lesson. 09:45–11:45.\par\textbf{Fluency routine:} R029, R053.\par\textbf{In class:} C07-U2; worked examples and applications; C07-E2 and E5. \par\textbf{Homework:} CQFD, exercise 21 p. 279, parts a and b. Write the condition of orthogonality to the normal, then check that the given point lies on the line. 10–20 min maximum; replaces the initial homework.\par\textbf{Note:} Expected mastery, to be adjusted from the work actually done and the diagnostics. CQFD solutions in the teacher booklet; return to it at the next lesson.
\Needspace{8\baselineskip}\section*{S30 -- 2027-01-12: Equation of a circle}
\textbf{Status:} lesson. 09:45–11:45.\par\textbf{Fluency routine:} R030, R054.\par\textbf{In class:} C07-U3, C07-U4; worked examples and applications; C07-E3 and E4.  MCQ 1 C07 (10 min, including correction). MCQ 2 C07 (10 min) and wrap-up.\par\textbf{Homework:} C07-E3, 10–20 min maximum.\par\textbf{Note:} Expected mastery, to be adjusted from the work actually done and the diagnostics.
\Needspace{8\baselineskip}\section*{S31 -- 2027-01-14: Slot reserved for mock exams}
\textbf{Status:} neutralised. 09:45–11:45.\par\textbf{Fluency routine:} none.\par\textbf{In class:} Mock exam 1 window (14–16 January): exact placement of mathematics to be confirmed.\par\textbf{Homework:} \par\textbf{Note:} Mock exam 1 window (14–16 January): exact placement of mathematics to be confirmed.
\Needspace{8\baselineskip}\section*{S32 -- 2027-01-19: Variable, event and distribution}
\textbf{Status:} lesson. 09:45–11:45.\par\textbf{Fluency routine:} R001, R056.\par\textbf{In class:} C08-U1; worked examples and applications; C08-E1. Diagnostic D: 10 min; investigation RECH: 15 min.\par\textbf{Homework:} C08-E1, 10–20 min maximum.\par\textbf{Note:} Expected mastery, to be adjusted from the work actually done and the diagnostics.
\Needspace{8\baselineskip}\section*{S33 -- 2027-01-21: Expected value and linearity}
\textbf{Status:} lesson. 09:45–11:45.\par\textbf{Fluency routine:} R002, R057.\par\textbf{In class:} C08-U2; worked examples and applications; C08-E3: fair stake only; variance deferred. \par\textbf{Homework:} CQFD, exercise 28 p. 306, all parts. Calculate the expected value from the distribution table; keep the probabilities associated with each value. 10–20 min maximum; replaces the initial homework.\par\textbf{Note:} Expected mastery, to be adjusted from the work actually done and the diagnostics. CQFD solutions in the teacher booklet; return to it at the next lesson.
\Needspace{8\baselineskip}\section*{S34 -- 2027-01-26: Variance, standard deviation and König–Huygens}
\textbf{Status:} lesson. 09:45–11:45.\par\textbf{Fluency routine:} R003, R058.\par\textbf{In class:} C08-U3; worked examples and applications; C08-E2.  MCQ 1 C08 (10 min, including correction).\par\textbf{Homework:} CQFD, exercise 31 p. 306, parts all parts; use the distribution from exercise 28. Calculate the variance then the standard deviation, respecting the rounding to the nearest tenth required by the textbook. 10–20 min maximum; replaces the initial homework.\par\textbf{Note:} Expected mastery, to be adjusted from the work actually done and the diagnostics. CQFD solutions in the teacher booklet; return to it at the next lesson.
\Needspace{8\baselineskip}\section*{S35 -- 2027-01-28: Transforming and comparing gains}
\textbf{Status:} lesson. 09:45–11:45.\par\textbf{Fluency routine:} R004, R059.\par\textbf{In class:} C08-U4, C08-U5; worked examples and applications; C08-E4 or E5 then sample workshop.  MCQ 2 C08 (10 min) and wrap-up.\par\textbf{Homework:} C08-E8 to revisit in a week, 10–20 min maximum.\par\textbf{Note:} Expected mastery, to be adjusted from the work actually done and the diagnostics.
\Needspace{8\baselineskip}\section*{S36 -- 2027-02-02: Defining a new function}
\textbf{Status:} lesson. 09:45–11:45.\par\textbf{Fluency routine:} R005, R054.\par\textbf{In class:} C09-U1; worked examples and applications; C09-AP1: tangent at 1. Diagnostic D: 10 min; investigation RECH: 15 min.\par\textbf{Homework:} C09-AP1, 10–20 min maximum.\par\textbf{Note:} Expected mastery, to be adjusted from the work actually done and the diagnostics.
\Needspace{8\baselineskip}\section*{S37 -- 2027-02-04: Algebraic properties and comparisons}
\textbf{Status:} lesson. 09:45–11:45.\par\textbf{Fluency routine:} R006, R055.\par\textbf{In class:} C09-U2; worked examples and applications; C09-E1 and E2. \par\textbf{Homework:} CQFD, exercise 18 p. 187, parts a, b and c. Simplify the products and quotients, citing the property of exponentials used. 10–20 min maximum; replaces the initial homework.\par\textbf{Note:} Expected mastery, to be adjusted from the work actually done and the diagnostics. CQFD solutions in the teacher booklet; return to it at the next lesson.
\Needspace{8\baselineskip}\section*{S38 -- 2027-02-09: Differentiating and studying}
\textbf{Status:} lesson. 09:45–11:45.\par\textbf{Fluency routine:} R007, R056.\par\textbf{In class:} C09-U3; worked examples and applications; C09-E3.  MCQ 1 C09 (10 min, including correction).\par\textbf{Homework:} CQFD, exercise 30 p. 187, parts a, b and c. Solve the three inequalities and write the solution sets as intervals. 10–20 min maximum; replaces the initial homework.\par\textbf{Note:} Expected mastery, to be adjusted from the work actually done and the diagnostics. CQFD solutions in the teacher booklet; return to it at the next lesson.
\Needspace{8\baselineskip}\section*{S39 -- 2027-02-11: From discrete to continuous}
\textbf{Status:} lesson. 09:45–11:45.\par\textbf{Fluency routine:} R008, R057.\par\textbf{In class:} C09-U4, C09-U5; worked examples and applications; C09-E4 and E5 as preferred.  MCQ 2 C09 (10 min) and wrap-up.\par\textbf{Homework:} C09-E8 to revisit in a week, 10–20 min maximum.\par\textbf{Note:} Expected mastery, to be adjusted from the work actually done and the diagnostics.
\Needspace{8\baselineskip}\section*{S40 -- 2027-03-02: Measuring an angle in radians}
\textbf{Status:} lesson. 09:45–11:45.\par\textbf{Fluency routine:} R009, R052.\par\textbf{In class:} C10-U1; worked examples and applications; C10-E1. Diagnostic D: 10 min; investigation RECH: 15 min.\par\textbf{Homework:} C10-E1, 10–20 min maximum.\par\textbf{Note:} Expected mastery, to be adjusted from the work actually done and the diagnostics.
\Needspace{8\baselineskip}\section*{S41 -- 2027-03-04: Wrapping and image of a real number}
\textbf{Status:} lesson. 09:45–11:45.\par\textbf{Fluency routine:} R010, R053.\par\textbf{In class:} C10-U2; worked examples and applications; C10-E2. \par\textbf{Homework:} C10-E2, 10–20 min maximum.\par\textbf{Note:} Expected mastery, to be adjusted from the work actually done and the diagnostics.
\Needspace{8\baselineskip}\section*{S42 -- 2027-03-09: Sine, cosine and exact values}
\textbf{Status:} lesson. 09:45–11:45.\par\textbf{Fluency routine:} R011, R054.\par\textbf{In class:} C10-U3; worked examples and applications; C10-E3 and E4.  MCQ 1 C10 (10 min, including correction).\par\textbf{Homework:} CQFD, exercise 39 p. 220, questions 1, 2, 3 and 5. Do the four questions on calculation and reading the circle; justify the answers in writing. 10–20 min maximum; replaces the initial homework.\par\textbf{Note:} Expected mastery, to be adjusted from the work actually done and the diagnostics. CQFD solutions in the teacher booklet; return to it at the next lesson.
\Needspace{8\baselineskip}\section*{S43 -- 2027-03-11: Associated angles and symmetries}
\textbf{Status:} lesson. 09:45–11:45.\par\textbf{Fluency routine:} R012, R055.\par\textbf{In class:} C10-U4, C10-U5; worked examples and applications; C10-E6 or E7.  MCQ 2 C10 (10 min) and wrap-up.\par\textbf{Homework:} CQFD, exercise 32 p. 220, parts a and b. Use the exact value given and the associated angles, without a calculator. 10–20 min maximum; replaces the initial homework.\par\textbf{Note:} Expected mastery, to be adjusted from the work actually done and the diagnostics. CQFD solutions in the teacher booklet; return to it at the next lesson.
\Needspace{8\baselineskip}\section*{S44 -- 2027-03-16: Differentiation and absolute value: remediation}
\textbf{Status:} lesson. 09:45–11:45.\par\textbf{Fluency routine:} R013, R050.\par\textbf{In class:} C04-E8; C06-E6; revisit the proofs for the square and the reciprocal.\par\textbf{Homework:} C06-E8, 10 min.\par\textbf{Note:} Expected mastery, to be adjusted from the work actually done and the diagnostics.
\Needspace{8\baselineskip}\section*{S45 -- 2027-03-18: Lists, sums and simulations}
\textbf{Status:} lesson. 09:45–11:45.\par\textbf{Fluency routine:} R014, R057.\par\textbf{In class:} Python workshop: suites.py and probabilites.py; C01-E7 and C08-E7, tasks chosen according to needs.\par\textbf{Homework:} Reread a program and write its output by hand, 15 min.\par\textbf{Note:} Expected mastery, to be adjusted from the work actually done and the diagnostics.
\Needspace{8\baselineskip}\section*{S46 -- 2027-03-23: Random variables: consolidating the models}
\textbf{Status:} lesson. 09:45–11:45.\par\textbf{Fluency routine:} R015, R064.\par\textbf{In class:} C08-E2 and E8; distribution, expected value and variance; protocol for simulating samples.\par\textbf{Homework:} C08-E5, 15 min.\par\textbf{Note:} Expected mastery, to be adjusted from the work actually done and the diagnostics.
\Needspace{8\baselineskip}\section*{S47 -- 2027-03-25: Preparing mock exam 2: methods}
\textbf{Status:} lesson. 09:45–11:45.\par\textbf{Fluency routine:} R016, R071.\par\textbf{In class:} EV3-T1: exercises, with an individual solution and a comparison of methods.\par\textbf{Homework:} Finish a designated exercise from paper EV3-T1, 20 min.\par\textbf{Note:} Expected mastery, to be adjusted from the work actually done and the diagnostics.
\Needspace{8\baselineskip}\section*{S48 -- 2027-03-30: Preparing mock exam 2: independence}
\textbf{Status:} lesson. 09:45–11:45.\par\textbf{Fluency routine:} R017, R031.\par\textbf{In class:} EV3-T2: selection of two exercises according to the diagnostic; verbalise the checks.\par\textbf{Homework:} Review the personal error sheet, 15 min.\par\textbf{Note:} Expected mastery, to be adjusted from the work actually done and the diagnostics.
\Needspace{8\baselineskip}\section*{S49 -- 2027-04-01: Slot reserved for mock exams}
\textbf{Status:} neutralised. 09:45–11:45.\par\textbf{Fluency routine:} none.\par\textbf{In class:} Mock exam 2 window (1–3 April): exact placement of mathematics to be confirmed.\par\textbf{Homework:} \par\textbf{Note:} Mock exam 2 window (1–3 April): exact placement of mathematics to be confirmed.
\Needspace{8\baselineskip}\section*{S50 -- 2027-04-06: Mock exam 2 feedback}
\textbf{Status:} lesson. 09:45–11:45.\par\textbf{Fluency routine:} R019, R045.\par\textbf{In class:} Correct the errors actually observed; papers EV3-A/B depending on the version sat.\par\textbf{Homework:} Redo an incorrect question without consulting the solutions, 15 min.\par\textbf{Note:} Expected mastery, to be adjusted from the work actually done and the diagnostics.
\Needspace{8\baselineskip}\section*{S51 -- 2027-04-08: Consolidation after mock exam 2}
\textbf{Status:} lesson. 09:45–11:45.\par\textbf{Fluency routine:} R020, R052.\par\textbf{In class:} C07-E2, E3 and C10-E6; two groups according to geometry errors.\par\textbf{Homework:} C07-E8, 10 min.\par\textbf{Note:} Expected mastery, to be adjusted from the work actually done and the diagnostics.
\Needspace{8\baselineskip}\section*{S52 -- 2027-04-27: Choosing a method in algebra}
\textbf{Status:} lesson. 09:45–11:45.\par\textbf{Fluency routine:} R021, R059.\par\textbf{In class:} C02-E3, E4 and E7\par\textbf{Homework:} Redo one targeted question, 10–15 min; no extra homework after a full paper.\par\textbf{Note:} Expected mastery, to be adjusted from the work actually done and the diagnostics.
\Needspace{8\baselineskip}\section*{S53 -- 2027-04-29: Probability: identifying the denominator}
\textbf{Status:} lesson. 09:45–11:45.\par\textbf{Fluency routine:} R022, R066.\par\textbf{In class:} C05-E1, E2 and E6\par\textbf{Homework:} Redo one targeted question, 10–15 min; no extra homework after a full paper.\par\textbf{Note:} Expected mastery, to be adjusted from the work actually done and the diagnostics.
\Needspace{8\baselineskip}\section*{S54 -- 2027-05-04: Differentiation: sign and domain}
\textbf{Status:} lesson. 09:45–11:45.\par\textbf{Fluency routine:} R023, R073.\par\textbf{In class:} C06-E1, E3 and E8\par\textbf{Homework:} Redo one targeted question, 10–15 min; no extra homework after a full paper.\par\textbf{Note:} Expected mastery, to be adjusted from the work actually done and the diagnostics.
\Needspace{8\baselineskip}\section*{S55 -- 2027-05-11: Geometry: choosing a framework}
\textbf{Status:} lesson. 09:45–11:45.\par\textbf{Fluency routine:} R024, R033.\par\textbf{In class:} C03-E3 and C07-E2\par\textbf{Homework:} Redo one targeted question, 10–15 min; no extra homework after a full paper.\par\textbf{Note:} Expected mastery, to be adjusted from the work actually done and the diagnostics.
\Needspace{8\baselineskip}\section*{S56 -- 2027-05-13: Exponential and sequences}
\textbf{Status:} lesson. 09:45–11:45.\par\textbf{Fluency routine:} R025, R040.\par\textbf{In class:} C09-E2, E3 and E8\par\textbf{Homework:} Redo one targeted question, 10–15 min; no extra homework after a full paper.\par\textbf{Note:} Expected mastery, to be adjusted from the work actually done and the diagnostics.
\Needspace{8\baselineskip}\section*{S57 -- 2027-05-18: EAM: timed practice T1}
\textbf{Status:} lesson. 09:45–11:45.\par\textbf{Fluency routine:} none.\par\textbf{In class:} EV4-T1, full paper 120 min\par\textbf{Homework:} Redo one targeted question, 10–15 min; no extra homework after a full paper.\par\textbf{Note:} Expected mastery, to be adjusted from the work actually done and the diagnostics.
\Needspace{8\baselineskip}\section*{S58 -- 2027-05-20: Practice T1 feedback}
\textbf{Status:} lesson. 09:45–11:45.\par\textbf{Fluency routine:} R027, R054.\par\textbf{In class:} EV4-T1: differentiated correction; redo two questions\par\textbf{Homework:} Redo one targeted question, 10–15 min; no extra homework after a full paper.\par\textbf{Note:} Expected mastery, to be adjusted from the work actually done and the diagnostics.
\Needspace{8\baselineskip}\section*{S59 -- 2027-05-25: Variance and statistics}
\textbf{Status:} lesson. 09:45–11:45.\par\textbf{Fluency routine:} R028, R061.\par\textbf{In class:} C08-E2 and E4; R017–R020 and R026–R028\par\textbf{Homework:} Redo one targeted question, 10–15 min; no extra homework after a full paper.\par\textbf{Note:} Expected mastery, to be adjusted from the work actually done and the diagnostics.
\Needspace{8\baselineskip}\section*{S60 -- 2027-05-27: Circle, radians and proofs}
\textbf{Status:} lesson. 09:45–11:45.\par\textbf{Fluency routine:} R029, R068.\par\textbf{In class:} C10-E2, E3 and E6; proofs of the exact values\par\textbf{Homework:} Redo one targeted question, 10–15 min; no extra homework after a full paper.\par\textbf{Note:} Expected mastery, to be adjusted from the work actually done and the diagnostics.
\Needspace{8\baselineskip}\section*{S61 -- 2027-06-01: EAM: timed practice T2}
\textbf{Status:} lesson. 09:45–11:45.\par\textbf{Fluency routine:} none.\par\textbf{In class:} EV4-T2, full paper 120 min\par\textbf{Homework:} Redo one targeted question, 10–15 min; no extra homework after a full paper.\par\textbf{Note:} Expected mastery, to be adjusted from the work actually done and the diagnostics.
\Needspace{8\baselineskip}\section*{S62 -- 2027-06-03: Practice T2 feedback}
\textbf{Status:} lesson. 09:45–11:45.\par\textbf{Fluency routine:} R078, R035.\par\textbf{In class:} EV4-T2: differentiated correction; individual plan\par\textbf{Homework:} Redo one targeted question, 10–15 min; no extra homework after a full paper.\par\textbf{Note:} Expected mastery, to be adjusted from the work actually done and the diagnostics.
\Needspace{8\baselineskip}\section*{S63 -- 2027-06-08: Consolidating proofs}
\textbf{Status:} lesson. 09:45–11:45.\par\textbf{Fluency routine:} R001, R042.\par\textbf{In class:} TRANS: logic; C01-U4, C03-U4 and C06-U2\par\textbf{Homework:} Redo one targeted question, 10–15 min; no extra homework after a full paper.\par\textbf{Note:} Expected mastery, to be adjusted from the work actually done and the diagnostics.
\Needspace{8\baselineskip}\section*{S64 -- 2027-06-10: Algebra-analysis synthesis}
\textbf{Status:} lesson. 09:45–11:45.\par\textbf{Fluency routine:} R002, R049.\par\textbf{In class:} C06-E5 and C09-E6\par\textbf{Homework:} Redo one targeted question, 10–15 min; no extra homework after a full paper.\par\textbf{Note:} Expected mastery, to be adjusted from the work actually done and the diagnostics.
\Needspace{8\baselineskip}\section*{S65 -- 2027-06-15: Fluency routines and reading data}
\textbf{Status:} lesson. 09:45–11:45.\par\textbf{Fluency routine:} R003, R056.\par\textbf{In class:} R001–R030, selection according to errors; R078\par\textbf{Homework:} Redo one targeted question, 10–15 min; no extra homework after a full paper.\par\textbf{Note:} Expected mastery, to be adjusted from the work actually done and the diagnostics.
\Needspace{8\baselineskip}\section*{S66 -- 2027-06-17: Final revisiting and exam technique}
\textbf{Status:} lesson. 09:45–11:45.\par\textbf{Fluency routine:} R004, R063.\par\textbf{In class:} Questions already missed; managing the 120 minutes, checking domains\par\textbf{Homework:} Redo one targeted question, 10–15 min; no extra homework after a full paper.\par\textbf{Note:} Expected mastery, to be adjusted from the work actually done and the diagnostics.
\Needspace{8\baselineskip}\section*{S67 -- 2027-06-22: Consolidation depending on whether lessons continue}
\textbf{Status:} conditional. 09:45–11:45.\par\textbf{Fluency routine:} R005, R070.\par\textbf{In class:} Targeted revisiting or transition to Terminale (Year 13): C06-E8, C09-E8, C10-E8, depending on attendance and exams.\par\textbf{Homework:} \par\textbf{Note:} Do not count this slot in the compulsory coverage before the EAM. Continuation of lessons to be confirmed.
\Needspace{8\baselineskip}\section*{S68 -- 2027-06-24: Consolidation depending on whether lessons continue}
\textbf{Status:} conditional. 09:45–11:45.\par\textbf{Fluency routine:} R006, R077.\par\textbf{In class:} Targeted revisiting or transition to Terminale (Year 13): C06-E8, C09-E8, C10-E8, depending on attendance and exams.\par\textbf{Homework:} \par\textbf{Note:} Do not count this slot in the compulsory coverage before the EAM. Continuation of lessons to be confirmed.
\Needspace{8\baselineskip}\section*{S69 -- 2027-06-29: Consolidation depending on whether lessons continue}
\textbf{Status:} conditional. 09:45–11:45.\par\textbf{Fluency routine:} R007, R037.\par\textbf{In class:} Targeted revisiting or transition to Terminale (Year 13): C06-E8, C09-E8, C10-E8, depending on attendance and exams.\par\textbf{Homework:} \par\textbf{Note:} Do not count this slot in the compulsory coverage before the EAM. Continuation of lessons to be confirmed.
\Needspace{8\baselineskip}\section*{S70 -- 2027-07-01: Consolidation depending on whether lessons continue}
\textbf{Status:} conditional. 09:45–11:45.\par\textbf{Fluency routine:} R008, R044.\par\textbf{In class:} Targeted revisiting or transition to Terminale (Year 13): C06-E8, C09-E8, C10-E8, depending on attendance and exams.\par\textbf{Homework:} \par\textbf{Note:} Do not count this slot in the compulsory coverage before the EAM. Continuation of lessons to be confirmed.
